% Options for packages loaded elsewhere
\PassOptionsToPackage{unicode}{hyperref}
\PassOptionsToPackage{hyphens}{url}
\documentclass[
]{article}
\usepackage{xcolor}
\usepackage[margin=1in]{geometry}
\usepackage{amsmath,amssymb}
\setcounter{secnumdepth}{-\maxdimen} % remove section numbering
\usepackage{iftex}
\ifPDFTeX
  \usepackage[T1]{fontenc}
  \usepackage[utf8]{inputenc}
  \usepackage{textcomp} % provide euro and other symbols
\else % if luatex or xetex
  \usepackage{unicode-math} % this also loads fontspec
  \defaultfontfeatures{Scale=MatchLowercase}
  \defaultfontfeatures[\rmfamily]{Ligatures=TeX,Scale=1}
\fi
\usepackage{lmodern}
\ifPDFTeX\else
  % xetex/luatex font selection
\fi
% Use upquote if available, for straight quotes in verbatim environments
\IfFileExists{upquote.sty}{\usepackage{upquote}}{}
\IfFileExists{microtype.sty}{% use microtype if available
  \usepackage[]{microtype}
  \UseMicrotypeSet[protrusion]{basicmath} % disable protrusion for tt fonts
}{}
\makeatletter
\@ifundefined{KOMAClassName}{% if non-KOMA class
  \IfFileExists{parskip.sty}{%
    \usepackage{parskip}
  }{% else
    \setlength{\parindent}{0pt}
    \setlength{\parskip}{6pt plus 2pt minus 1pt}}
}{% if KOMA class
  \KOMAoptions{parskip=half}}
\makeatother
\usepackage{longtable,booktabs,array}
\newcounter{none} % for unnumbered tables
\usepackage{calc} % for calculating minipage widths
% Correct order of tables after \paragraph or \subparagraph
\usepackage{etoolbox}
\makeatletter
\patchcmd\longtable{\par}{\if@noskipsec\mbox{}\fi\par}{}{}
\makeatother
% Allow footnotes in longtable head/foot
\IfFileExists{footnotehyper.sty}{\usepackage{footnotehyper}}{\usepackage{footnote}}
\makesavenoteenv{longtable}
\setlength{\emergencystretch}{3em} % prevent overfull lines
\providecommand{\tightlist}{%
  \setlength{\itemsep}{0pt}\setlength{\parskip}{0pt}}
\usepackage{bookmark}
\IfFileExists{xurl.sty}{\usepackage{xurl}}{} % add URL line breaks if available
\urlstyle{same}
\hypersetup{
  hidelinks,
  pdfcreator={LaTeX via pandoc}}

\author{}
\date{}

\begin{document}

{
\setcounter{tocdepth}{3}
\tableofcontents
}
\section{A Systematic Approach for LLM-Based Functional Software Test
Case Generation and Autonomous
Execution}\label{a-systematic-approach-for-llm-based-functional-software-test-case-generation-and-autonomous-execution}

\textbf{A reproduction-and-extension study.} Formatted after the
IntelliTest capstone thesis (\emph{A Systematic Approach for LLM-Based
Functional Software Test Case Generation}) and the prior published paper
(\emph{A Business Process-Centric Approach for LLM-Driven Functional
Test Generation}).

\begin{quote}
\textbf{Scope and honesty statement.} The \textbf{generation} pipeline
(§III-A) is \textbf{prior published work} and is reproduced here, not
claimed as novel. The contributions of \emph{this} document are: (1) a
faithful \textbf{reproduction on a new multi-model backend} (BytePlus:
\texttt{deepseek-v4-flash}, \texttt{deepseek-v4-pro},
\texttt{dola-seed-2-1-turbo}) with complete per-call token/cost/latency
logging; (2) an \textbf{end-to-end integration} of generation with the
Test Execution and Adaptation Agent (TEAA) against a live,
locally-deployed system under test (FullTeaching); and (3) a
\textbf{transparent economic and reliability analysis} that discloses,
rather than hides, discrepancies between metrics computed over different
test populations. Every quantitative claim traces to a logged artifact
under \texttt{research\_project/}.
\end{quote}

\begin{center}\rule{0.5\linewidth}{0.5pt}\end{center}

\subsection{Definitions and Acronyms}\label{definitions-and-acronyms}

{\def\LTcaptype{none} % do not increment counter
\begin{longtable}[]{@{}ll@{}}
\toprule\noalign{}
Acronym & Definition \\
\midrule\noalign{}
\endhead
\bottomrule\noalign{}
\endlastfoot
API & Application Programming Interface \\
BFS & Breadth-First Search \\
BVA & Boundary Value Analysis \\
BVC & Boundary Value Coverage \\
BP & Business Process \\
CoT & Chain-of-Thought \\
DOM & Document Object Model \\
E2E & End-to-End \\
EP & Equivalence Partitioning \\
LLM & Large Language Model \\
LTM & Long-Term Memory \\
SRS & Software Requirements Specification \\
STC & State Transition Coverage \\
STM & Short-Term Memory \\
SUT & System Under Test \\
TEAA & Test Execution and Adaptation Agent \\
UAT & User Acceptance Testing \\
VBER & Verified BVA Execution Rate \\
VLM & Vision-Language Model \\
\end{longtable}
}

\begin{center}\rule{0.5\linewidth}{0.5pt}\end{center}

\subsection{Abstract}\label{abstract}

Validating end-to-end business processes in software engineering remains
a critical bottleneck because of the \emph{semantic gap} between
ambiguous textual requirements and the precision required for executable
test scenarios. Prior published work --- the business-process-centric
IntelliTest \textbf{generation} pipeline --- demonstrated that combining
the semantic reasoning of Large Language Models (LLMs) with
deterministic Specification-Based Testing (Equivalence Partitioning,
Boundary Value Analysis, and Pairwise Testing), together with a
graph-based navigational model, produces mathematically complete,
high-coverage functional test suites while mitigating the ``context
amnesia'' of purely generative models. That work stops at test
\emph{generation}.

This paper reproduces the generation pipeline on a new,
provider-agnostic multi-model backend (BytePlus, OpenAI-compatible) and
instruments it so that \textbf{every} model call records its prompt,
response, prompt and completion token counts, monetary cost, and
latency. We then integrate the generation output with the \textbf{Test
Execution and Adaptation Agent (TEAA)}, an engine that translates the
structured, ISO/IEC/IEEE 29119-3-style test specifications into
resilient interactions on a live web application. We deploy the system
under test (FullTeaching, an Angular + Spring + OpenVidu
learning-management platform) locally via Docker and confirm the full
generation→execution loop.

Empirically, a full generation run over the 31-page FullTeaching SRS
consumes \textbf{2,450,312 tokens across 69 model calls
(deepseek-v4-flash)} or approximately \textbf{2.6M tokens across 90
calls (deepseek-v4-pro)}, producing 63 ISO-style test cases and 100+
mutation test cases in roughly 50--63 minutes. The two models produce
comparable artifact volumes, but the ``pro'' model costs approximately
\textbf{3.5× more} per suite under current (placeholder) pricing. On the
execution side, we verify from the authors' artifacts a \textbf{42.2\%
all-flows success rate} and an \textbf{81.3\% valid-action rate} over
166 executed test cases and 1,756 browser actions. We explicitly
disclose that this all-flows figure differs from the thesis's headline
89.2\% normal-flow figure because the two are computed over different
populations, and that the observed exception-flow success (5.9\%)
contradicts the thesis's ``\textgreater80\%'' exception claim --- we
report both and treat exception handling as a documented limitation
rather than a strength.

\textbf{Keywords:} Large Language Models, software testing, automated
test generation, agentic execution, functional testing, boundary value
analysis, pairwise testing.

\begin{center}\rule{0.5\linewidth}{0.5pt}\end{center}

\subsection{I. Introduction}\label{i.-introduction}

\subsubsection{I-A. The system-testing
bottleneck}\label{i-a.-the-system-testing-bottleneck}

The Software Development Lifecycle treats testing not as a final phase
but as a parallel assurance process. Within it, \textbf{System Testing}
is the definitive quality gate: it verifies that software complies with
the business requirements captured in the Software Requirements
Specification (SRS). Achieving high-fidelity assurance in this phase is
cognitively expensive. Testers must (1) decompose the documentation into
discrete components --- functional logic, UI designs --- and map the web
of interactions and state transitions between them; and (2)
strategically apply formal specification-based techniques to narrow an
effectively infinite input space into a finite, high-coverage test
suite.

This burden is twofold. The \textbf{Test Analysis and Design} phase
demands substantial cognitive load to abstract business logic into
executable scenarios. The \textbf{Test Implementation} phase --- whether
by manual execution or by writing automation scripts --- consumes large
amounts of time. The high cost of maintaining this rigor commonly
produces a \textbf{happy-path bias}: testers prioritize valid workflows
and neglect the edge cases where defects concentrate, in order to meet
delivery deadlines.

\subsubsection{I-B. Why naive LLM use is
insufficient}\label{i-b.-why-naive-llm-use-is-insufficient}

LLMs promise to automate this domain, but current approaches exhibit two
architectural deficiencies that this project takes as its motivating
premises:

\begin{enumerate}
\def\labelenumi{\arabic{enumi}.}
\item
  \textbf{The rigor gap (coverage reliability).} The stochastic nature
  of LLMs leads to hallucination and precludes the systematic
  application of formal techniques. Consider a screen with \texttt{n}
  input parameters, each with a domain of representative values
  \texttt{D\_i} after Equivalence Partitioning. The exhaustive test
  space is the Cartesian product
  \texttt{\textbar{}Ω\_data\textbar{}\ =\ Π\textbar{}D\_i\textbar{}}. A
  registration form with \texttt{n=10} fields and 4 partitions each
  yields \texttt{4\^{}10\ ≈\ 1,048,576} combinations. A purely
  inferential model produces a stochastic subset \texttt{T̂\ ⊂\ Ω\_data}
  with \texttt{\textbar{}T̂\textbar{}\ ≪\ \textbar{}Ω\_data\textbar{}},
  heavily biased toward the happy path. Combinatorial optimization
  (Pairwise Testing) is required to reduce this space while preserving
  interaction coverage.
\item
  \textbf{Context amnesia (contextual scope).} LLMs operate as localized
  inference engines and tend to generate tests for functions in
  isolation, failing to synthesize the stateful, multi-step business
  workflows that validate end-to-end logic. Formally, modeling the
  application as a state-transition graph \texttt{G\ =\ (V,\ E)}, a
  business process is a directed path
  \texttt{π\ =\ ⟨s\_0\ →a\_1\ s\_1\ →a\_2\ …\ →a\_k\ s\_k⟩}, and
  validating state \texttt{s\_k} often depends on latent variables
  produced at earlier steps. Purely generative models lose this
  continuity.
\end{enumerate}

A parallel deficiency appears in \textbf{execution}: autonomous agents
often navigate without a grounded script (behaving like an ``unguided
mouse in a maze'') and, crucially, lack a definitive \textbf{oracle} ---
an awareness of correct behavior derived from requirements --- so they
cannot reliably decide whether the application is functioning correctly.

\subsubsection{I-C. Contributions of this
document}\label{i-c.-contributions-of-this-document}

The prior published IntelliTest generation pipeline addresses the rigor
gap and context amnesia for the \emph{design} phase, using LLMs strictly
for requirement comprehension while delegating completeness to
deterministic algorithms (EP, BVA, Pairwise) and a graph-based
navigational model. Because that work is already published, this
document does not re-claim it. Instead, the contributions here are
deliberately scoped to what is new:

\begin{enumerate}
\def\labelenumi{\arabic{enumi}.}
\tightlist
\item
  \textbf{Reproduction on a new multi-model backend with full
  accounting.} We port the pipeline to BytePlus (an OpenAI-compatible
  endpoint) and run it on three models, logging every call's tokens,
  cost, and latency. This establishes that the method is not tied to a
  single model family and yields the per-model cost profile the original
  did not report.
\item
  \textbf{End-to-end generation→execution integration.} We connect the
  generated ISO-style specifications to the TEAA execution agent and
  confirm the loop against a \textbf{live}, locally-deployed
  FullTeaching SUT.
\item
  \textbf{Transparent economics and reliability accounting.} We report
  measured token/cost figures and, per the supervisor's explicit
  instruction to \emph{not hide any data}, we disclose where measured
  execution metrics diverge from previously reported headline numbers,
  including a contradiction in the exception-flow claim.
\end{enumerate}

The remainder is organized as: §II literature review (an 82-paper PRISMA
synthesis, with Yun Lin's WebTestPilot as the closest execution-phase
prior art); §III methodology (generation as referenced prior work;
execution/TEAA as the integration focus); §IV experiment; §V discussion
including economic feasibility and threats to validity; §VI conclusion.

\begin{center}\rule{0.5\linewidth}{0.5pt}\end{center}

\subsection{II. Literature Review}\label{ii.-literature-review}

\subsubsection{II-A. Review protocol
(PRISMA)}\label{ii-a.-review-protocol-prisma}

To situate this work, we conducted a PRISMA-style systematic review.
Search strings combined four keyword groups --- \textbf{domain} \{system
testing, functional testing, UAT, acceptance testing\} ×
\textbf{technique} \{LLM, large language model, AI agent, autonomous,
multi-agent\} × \textbf{target} \{web application, GUI, interface, app\}
× \textbf{activity} \{test case generation, test execution, test
automation\} --- augmented with adjacent subfields (unit/integration
test generation, oracle generation, self-healing automation,
DOM/accessibility abstraction, RAG for testing, VLM GUI testing, REST
API testing, BDD/Gherkin, flaky-test detection, requirements
traceability, combinatorial testing). Approximately 30 distinct queries
were executed against arXiv, IEEE Xplore, the ACM Digital Library, and
venue proceedings. After title/abstract screening and de-duplication by
normalized title and arXiv identifier, \textbf{82 papers} were included.
The machine-readable corpus (11-column schema: title, authors, year,
venue, link, quartile, summary, problem solved, methodology, type,
original link) is provided as \texttt{related\_articles.xlsx} across
three sheets (Systematic Review = 82; New Research 2025--26 = 19; All
References = 101 rows). Uncertain author/venue/quartile fields are
explicitly marked \texttt{{[}estimated{]}}; all arXiv identifiers are
real.

\subsubsection{II-B. Taxonomy of dynamic testing
techniques}\label{ii-b.-taxonomy-of-dynamic-testing-techniques}

We frame the work within the standard taxonomy. \textbf{Structure-based
(white-box)} techniques (statement, decision coverage) verify the
``how'' of implementation but are insufficient for functional
validation. \textbf{Experience-based} techniques (error guessing,
exploratory testing) are valuable but subjective and hard to automate
reproducibly. The focus of this project is \textbf{Specification-Based
Testing (black-box)}, which derives tests strictly from external
documentation (SRS, use cases, user stories). Within it:
\textbf{Equivalence Partitioning (EP)} and \textbf{Boundary Value
Analysis (BVA)} are the input-validation pillars; \textbf{Pairwise
Testing} mitigates combinatorial explosion by guaranteeing that every
pair of parameter values is covered at least once; \textbf{Decision
Table} and \textbf{State Transition} testing address logic and lifecycle
dynamics; and \textbf{Use Case Testing} links functions into coherent
workflows.

\subsubsection{II-C. Two paradigms and the gap between
them}\label{ii-c.-two-paradigms-and-the-gap-between-them}

The 82-paper corpus separates into two paradigms:

\begin{itemize}
\tightlist
\item
  \textbf{Top-down, document/requirement-driven GENERATION.} The prior
  IntelliTest work sits here. Neighbours: ChatTester and ChatUniTest
  (unit-test generation with adaptive focal context and
  generation-validation- repair loops); high-level/requirements test
  generation (arXiv:2503.17998, 2510.03641); UAT tooling (AutoUAT + Test
  Flow, arXiv:2504.07244; AToMIC, arXiv:2510.18861); and multi-agent
  generators with oracle panels (CANDOR, arXiv:2506.02943; TestAgent,
  arXiv:2607.09101).
\item
  \textbf{Bottom-up, AGENTIC EXECUTION.} DroidAgent (intent-driven
  mobile GUI testing with long/short-term memory), GPTDroid
  (functionality-aware GUI Q\&A), WebVoyager (multimodal end-to-end web
  agent, 59.1\% task success), and Autonomous Test Agents
  (PinATA/SeeAct-ATA, arXiv:2504.01495).
\end{itemize}

\textbf{Three synthesis claims} carry into this paper:

\begin{enumerate}
\def\labelenumi{\arabic{enumi}.}
\tightlist
\item
  \textbf{Two paradigms, one gap.} No single 2025--2026 work integrates
  \emph{business-process-grounded generation} with an \emph{autonomous
  UAT execution agent} that issues reliable verdicts, while
  simultaneously addressing soundness/consistency and oracle
  reliability. The prior published work covers generation only; this
  document's contribution is the integration.
\item
  \textbf{The rigor gap is measurable.} Generation is evaluated with
  reproducible metrics (coverage, mutation, assertion correctness).
  Execution/agentic work relies on ad-hoc task-success rates with no
  shared oracle benchmark --- WebVoyager reports 59.1\% task success;
  PinATA reports \textasciitilde60\% correct verdicts (up to 94\%
  specificity). This is precisely why our execution results (§IV-C) are
  reported against multiple denominators.
\item
  \textbf{Context amnesia is shared.} ChatUniTest's token-budget focal
  context and DroidAgent's explicit memory modules both exist to fight
  the same amnesia the generation pipeline threads business-process
  context through. Our TEAA carries this forward with short-term and
  long-term (Golden Trace) memory.
\end{enumerate}

\subsubsection{II-D. Closest prior art for the execution
phase}\label{ii-d.-closest-prior-art-for-the-execution-phase}

Two works are the primary anchors for the execution half:

\begin{itemize}
\item
  \textbf{WebTestPilot} --- Xiwen Teoh, \textbf{Yun Lin} (Associate
  Professor \& Deputy Head of the CS Department, Shanghai Jiao Tong
  University; formerly NUS, PAT group of Prof.~Jin Song Dong), Duc-Minh
  Nguyen, Ruofei Ren, Wenjie Zhang, Jin Song Dong. \textbf{FSE 2026}
  (arXiv:2602.11724). This is the single closest concurrent competitor
  to the TEAA. It is a \textbf{neurosymbolic} LLM agent that (1) detects
  and \textbf{symbolizes} critical GUI elements into variables, and (2)
  translates an NL specification into a sequence of steps, each equipped
  with \textbf{inferred pre- and post-condition oracles} over those
  symbols, capturing \textbf{data, temporal, and causal} dependencies.
  It directly attacks the two problems the thesis also names --- the
  \emph{implicit oracle inference} problem (the agent must be its own
  oracle) and the \emph{probabilistic inference} problem (distinguishing
  a hallucination from a real bug). On a bug-injected benchmark it
  reports \textbf{99\% task completion, 96\% precision, and 96\% recall}
  (+70 precision / +27 recall over the best baseline).
  \emph{Differentiation:} our TEAA grounds execution in the
  \textbf{business-process graph inherited from the generation half},
  and uses \textbf{Semantic DOM abstraction + Golden-Trace memory + a
  multi-modal (DOM diff + screenshot) oracle + trace mutation}, whereas
  WebTestPilot centers on symbolization and inferred symbolic
  pre/post-condition oracles. WebTestPilot is the recommended primary
  comparison for any future head-to-head execution evaluation.

  Yun Lin's group produces several other directly relevant 2025--2026
  works: \emph{Compiling Large Multi-Modal Requirement Documents into
  Runnable Software Systems} (ISSTA 2026, arXiv:2602.13723);
  \emph{Generating Project-Specific Test Cases with Requirement
  Validation Intention} (ISSTA 2026); \emph{Generalizing Test Cases for
  Comprehensive Test Scenario Coverage} (TestGeneralizer, FSE 2026,
  arXiv:2604.21771); and design-to- action agents for GUI testing (ISSTA
  2025).
\item
  \textbf{Are Autonomous Web Agents Good Testers? (PinATA / SeeAct-ATA)}
  --- arXiv:2504.01495. Establishes that autonomous test agents can
  execute NL test cases and return pass/fail verdicts, but at limited
  reliability (\textasciitilde60\% correct verdict). This is the
  empirical baseline for ``can agents judge correctness at all,'' and it
  contextualizes our 42.2\% all-flows figure.
\end{itemize}

\subsubsection{II-E. Newly-emerged concerns (2025--2026) the paper must
acknowledge}\label{ii-e.-newly-emerged-concerns-20252026-the-paper-must-acknowledge}

The re-research track surfaced several themes that postdate the prior
published work and shape the execution half's threats-to-validity:

\begin{itemize}
\tightlist
\item
  \textbf{NL-test soundness and consistency} (arXiv:2509.19136): NL
  tests can be \emph{unsound} (false failures from ambiguity) and
  \emph{inconsistent} (nondeterministic re-runs). This must appear as a
  threat to validity.
\item
  \textbf{Self-healing via DOM/accessibility trees} (arXiv:2603.20358):
  sub-second selector re-discovery at zero ongoing API cost, scaling to
  300+ tests.
\item
  \textbf{Bounded autonomous repair} (arXiv:2605.01471):
  \textasciitilde70\% repair convergence, mean 4.4 iterations ---
  realistic limits on self-correction during execution.
\item
  \textbf{DOM pruning at scale} (Prune4Web, arXiv:2511.21398):
  programmatic pruning of 10k--100k-token DOMs, the same bottleneck the
  TEAA's Semantic DOM Abstraction addresses.
\item
  \textbf{Accessibility-tree prompt injection} (arXiv:2507.14799): a
  genuine security threat model for DOM/ accessibility-tree-driven
  agents.
\item
  \textbf{Requirement→oracle alignment on real bugs} (arXiv:2607.10277):
  LLM-generated oracles align more with requirements than with the SUT
  and show high variance --- central to autonomous-verdict reliability.
\end{itemize}

\begin{center}\rule{0.5\linewidth}{0.5pt}\end{center}

\subsection{III. Methodology}\label{iii.-methodology}

\subsubsection{III-A. Test Case Design Phase (PRIOR PUBLISHED WORK ---
reproduced, not
claimed)}\label{iii-a.-test-case-design-phase-prior-published-work-reproduced-not-claimed}

The generation pipeline is the previously published IntelliTest design
phase. It is included here in detail only to ground the execution
integration and to document the faithful reproduction on BytePlus. It
runs eight stages end-to-end
(\texttt{intellitest/main.py::run\_full\_pipeline}), each funneling
through a single provider abstraction
(\texttt{LLMProvider.generate\_structured}) that returns a validated
Pydantic object.

\textbf{Stage 1 --- Screen extraction.} Extract the set of system
screens \texttt{S} from the SRS → \texttt{general/screens.json}. On the
FullTeaching SRS this yielded 4 screens (Landing, Dashboard, Course
Details, Video Session).

\textbf{Stage 2 --- Business Process Detector.} Extract ordered,
multi-step business processes. A \textbf{grounding constraint} prevents
hallucinated navigation: for every step \texttt{st} in every business
process \texttt{b}, \texttt{\{cs(st),\ ns(st)\}\ ⊆\ S\_names} --- the
current and next screen of each step must be members of the
pre-identified screen set. On FullTeaching this produced 6 business
processes (\texttt{business\_process\_1..6.json}).

\textbf{Stage 3 --- Screen Variable Detector.} For each screen, an LLM
function \texttt{F\_var} builds a behavioral model
\texttt{M\_var(S\_i)\ =\ F\_var(P\_var(B(S\_i)),\ C)} (eq. 3.6), where
\texttt{B(S\_i)} is the subset of business processes touching screen
\texttt{S\_i} and \texttt{C} is the document context. The prompt
introduces a \textbf{hierarchical layer structure}
\texttt{L\_i\ =\ \{l\_0,\ l\_1,\ …,\ l\_m\}} (a base layer of
always-visible elements; each subsequent layer gated by a trigger in the
previous), organizes elements into \textbf{element groups} \texttt{g},
and instructs the model to act as a Test Analyst applying BVA,
decision-table, and state-machine reasoning. A \textbf{Strict Abstract
Value Constraint} forces abstract data categories (e.g.,
\texttt{"valid\_email\_format"}) rather than concrete instances. For
each parameter \texttt{p} it yields a domain
\texttt{TD\_p\ =\ (V\_valid,\ V\_invalid)}. Each group is linked to its
business function:
\texttt{∀g\ ∈\ M\_var(S\_i):\ β(g)\ ∈\ \{b.name\ \textbar{}\ b\ ∈\ B(S\_i)\}}
(eq. 3.7).

\textbf{Stage 4 --- Screen Graph Detector.} Build the directed
navigational graph \texttt{G\_screen\ =\ (V,\ E)} with \texttt{V\ =\ S}.
For each screen, a CoT-guided function \texttt{F\_graph} emits the
outgoing edge subset \texttt{E\_i\ =\ F\_graph(P\_graph,\ C,\ s\_i)}
(eq. 3.5). Each edge is a tuple
\texttt{e\ =\ (s\_source,\ s\_target,\ a,\ el,\ K,\ k\_imp)} (action,
element, keyword set, most-critical keyword). A thread-safe writer
atomically unions the subsets: \texttt{E\ =\ ∪\ E\_i}.

\textbf{Stage 5 --- Path Processor.} For each step \texttt{st\_k} of
business process \texttt{b\_j}, compute the navigation path via BFS:
\texttt{π\_k\ =\ BFS(G\_screen,\ s\_start(π\_k),\ cs(st\_k))} (eq.
3.11), where the start screen carries continuity ---
\texttt{s\_start(π\_k)\ =\ cs(st\_1)} if \texttt{k=1}, else
\texttt{ns(st\_\{k-1\})} (eq. 3.10).

\textbf{Stage 6 --- Pairwise Generator.} For each element group's
collected parameter set \texttt{P\_gk} (gathered recursively across the
hierarchical layers), apply All-Pairs:
\texttt{Comb(g\_k)\ =\ Pairwise(TD\_p1,\ …,\ TD\_pm)} (eq. 3.8). The
expected outcome of a combination is deterministic:
\texttt{outcome(c)\ =\ False\ if\ c\ ∩\ V\_invalid\ ≠\ ∅\ else\ True}
(eq. 3.9).

\textbf{Stage 7 --- Test Case Generator.} For each
\texttt{(st\_k,\ c\ ∈\ Comb(g\_k))}, synthesize one test case
\texttt{tc\_k\ =\ F\_gen(P\_gen(st\_k,\ π\_k,\ c,\ H\_\{k-1\},\ D\_input),\ C)}
(eq. 3.12), where \texttt{H\_\{k-1\}} is the aggregated historical
context of prior successful test cases in the same flow --- enabling
\textbf{stateful testing} (e.g., reusing a newly created account to log
in) --- and \texttt{D\_input} are injected constants (e.g., admin
credentials). Output conforms to \textbf{ISO/IEC/IEEE 29119-3:2013}. A
\textbf{Test Case Retry} step re-attempts failed generations.

\textbf{Stage 8 --- Mutation / Exception Generator.} Synthesize
exception/negative cases by mutating tested rules:
\texttt{tc\textquotesingle{}\_k\ =\ F\_except(P\_except(st\_k,\ r),\ C)}
(eq. 3.13).

\textbf{Reproduction defect fix.} The BytePlus (OpenAI-compatible)
provider path originally set \textbf{no \texttt{max\_tokens}}, so the
large business-process JSON was truncated mid-string (``EOF while
parsing a string''), which silently zeroed every downstream stage (0
business processes → 0 paths → 0 pairwise → 0 test cases). The fix adds
a \texttt{max\_tokens} bound (env \texttt{DEEPSEEK\_MAX\_TOKENS},
default 16384) and per-call usage logging. This is the only functional
change; prompts and control flow are unchanged.

\subsubsection{III-B. Test Execution Phase --- TEAA (integration focus
of this
work)}\label{iii-b.-test-execution-phase-teaa-integration-focus-of-this-work}

The \textbf{Test Execution and Adaptation Agent} consumes the generated
ISO-style specifications and drives a live browser via Selenium. It is
implemented as \texttt{TestAgent} in \texttt{CustomLLMDroid/Agent.py}
(positive/negative Golden-Path pipeline) and
\texttt{CustomLLMDroid/Mutation\_Agent.py} (structural exception flows),
with LLM access in \texttt{LLMCaller.py}, prompts in \texttt{Prompt.py},
and JSON schemas in \texttt{Schemas.py}. The following table maps each
thesis concept (§3.2, eqs 3.14--3.20) to its concrete code realization
(verified by source inspection):

{\def\LTcaptype{none} % do not increment counter
\begin{longtable}[]{@{}
  >{\raggedright\arraybackslash}p{(\linewidth - 6\tabcolsep) * \real{0.2500}}
  >{\raggedright\arraybackslash}p{(\linewidth - 6\tabcolsep) * \real{0.2500}}
  >{\raggedright\arraybackslash}p{(\linewidth - 6\tabcolsep) * \real{0.2500}}
  >{\raggedright\arraybackslash}p{(\linewidth - 6\tabcolsep) * \real{0.2500}}@{}}
\toprule\noalign{}
\begin{minipage}[b]{\linewidth}\raggedright
Thesis concept
\end{minipage} & \begin{minipage}[b]{\linewidth}\raggedright
Formula
\end{minipage} & \begin{minipage}[b]{\linewidth}\raggedright
Code
\end{minipage} & \begin{minipage}[b]{\linewidth}\raggedright
Realization
\end{minipage} \\
\midrule\noalign{}
\endhead
\bottomrule\noalign{}
\endlastfoot
\textbf{Golden-Path filter}
\texttt{G\_k\ =\ \{tc\ :\ ExpectedOutcome\ =\ True\}} & 3.14 &
\texttt{Agent.walkthrough()} & Pops the positive-outcome test case as
\texttt{success\_tc} before running negatives. \\
\textbf{Semantic DOM Abstraction} \texttt{Φ:\ H\_t\ →\ S\_t}, keeping
visible interactive nodes & 3.15 &
\texttt{extract\_interactable\_elements()} + \texttt{is\_visible()} &
BeautifulSoup prune of
\texttt{display:none/visibility:hidden/opacity:0/aria-hidden} and hidden
ancestors (visual visibility \texttt{v(e)=1}); keep clickable/input/form
tags and elements with \texttt{role/on*/tabindex} (\texttt{A\_ctrl});
emit compact JSON (thesis claims \textasciitilde95\% token
reduction). \\
\textbf{Strategy Function} \texttt{a\_t\ =\ σ(S\_t,\ M\_STM,\ D\_input)}
& 3.16 & \texttt{evaluate\_next\_step()} + \texttt{EVALUATE\_NEXT\_STEP}
prompt/schema & Single-action decision loop; input = abstracted elements
\texttt{S\_t}, short-term memory \texttt{M\_STM} (execute log),
test-case fields, assumption + true \texttt{D\_input}; output =
\texttt{\{action,\ selector\_type,\ selector\_value,\ input\_value,\ reasoning,\ done,\ web\_fail,\ llm\_fail\}}. \\
\textbf{Termination} \texttt{σ(...)\ =\ "done"\ ⇒\ terminate} & 3.17 &
\texttt{while\ not\ done:} + 600s watchdog & Watchdog resets
driver/session to break infinite loops. \\
\textbf{Golden Trace / LTM} \texttt{τ\_verified\ →\ M\_LTM} &
(3.16--3.17) & \texttt{long\_term\_mem},
\texttt{success\_happy\_path.json},
\texttt{evaluate\_past\_success\_tc()} & On success, the error-free
execute log is persisted; a boosting step replays matching past traces
to bypass expensive re-inference. \\
\textbf{Multi-modal Oracle}
\texttt{Verdict\ =\ PASS\ iff\ ValidTech\ ∧\ Match(ΔS,\ E\_exp)} & 3.18
& \texttt{evaluate\_test\_execution()} + \texttt{get\_html\_diff()} +
images via \texttt{generate\_json\_with\_images()} & Captures
before/after screenshots (base64) + unified HTML diff (≤5000 chars) +
visible text; returns
\texttt{\{valid\_action,\ result,\ llm\_fail,\ web\_fail,\ reasoning,\ assertion\_selector\}}
--- a two-channel structural+visual verdict. \\
\textbf{Trace Mutation} \texttt{τ\_neg\ =\ Ψ(τ\_verified,\ δ)} & 3.19 &
\texttt{generate\_false\_testcase()}; structural in
\texttt{Mutation\_Agent.py} & Preserves navigation, perturbs input data
(invalid/empty/malformed); oracle logic reversed (expect
error/blocking). \\
\textbf{Deterministic Serialization}
\texttt{S\_exec\ =\ Γ(τ\_verified,\ A)} replayed by \texttt{R\_det} &
3.20 & \texttt{Reuseable\_Module.ReusableTestRunner}, \texttt{rerun.py}
& LLM-free Selenium replay of saved actions + assertions, with 12
assertion primitives; \texttt{C\_script\ ≪\ C\_LLM} --- no model call at
replay. \\
\end{longtable}
}

The key integration property is that \textbf{the generation half feeds
the execution half}: the business-process-grounded, ISO-style
specification becomes the TEAA's script, directly addressing the
``unguided explorer'' and ``missing oracle'' deficiencies of prior
agentic testing.

\begin{center}\rule{0.5\linewidth}{0.5pt}\end{center}

\subsection{IV. Experiment}\label{iv.-experiment}

\subsubsection{IV-A. Experimental setup}\label{iv-a.-experimental-setup}

\textbf{Case studies.} 1. \textbf{FullTeaching} --- an open-source
learning-management platform (Angular front end, Spring back end,
OpenVidu WebRTC), selected because of its rich dynamic interface and
complex state transitions. It is the live SUT for the execution phase.
Deployed locally via Docker Compose (three containers: the app, MySQL,
and the OpenVidu KMS server) and confirmed reachable at
\texttt{https://localhost:5000/} (HTTP 200; HTTPS only --- plain HTTP
returns 400). Default accounts: \texttt{teacher@gmail.com/pass},
\texttt{student1@gmail.com/pass}, \texttt{student2@gmail.com/pass}. Its
SRS (\texttt{full-teaching-system-design.pdf}, 31 pages, 60,717
characters extracted, tabulated functional requirements FR-001\ldots)
drives generation; the pipeline detected 4 screens and 6 business
processes. 2. \textbf{RetailOnboardPro / 4LV (limited sample).} The
available artifact \texttt{srs\_extracted\_4lv.json} is a
\textbf{2-use-case sample} (Submit Training Request; Approve Training
Request), \textbf{not} the full RetailOnboardPro SRS from the original
thesis (which has 10 use cases, 8 screens, 75 business rules). It was
rendered to PDF (\texttt{make\_4lv\_pdf.py}) and used to characterize
behavior on a minimal input. All 4LV results are reported explicitly as
a limited sample.

\textbf{Models (BytePlus, OpenAI-compatible,
\texttt{https://ark.ap-southeast.bytepluses.com/api/v3}).}
\texttt{deepseek-v4-flash-ga-260731} (flash),
\texttt{deepseek-v4-pro-ga-260813} (pro),
\texttt{dola-seed-2-1-turbo-260628} (dola). All three were verified to
respond and to report usage tokens. flash and pro were run through full
pipelines; dola was connectivity-validated but not run through a full
pipeline (a documented limitation).

\textbf{Baselines (generation).} Three faithful reimplementations,
confirmed by a code audit to be verbatim-prompt, control-flow-accurate
ports of their source methods: - \textbf{Baseline 1 (Augusto et al.,
ICTSS 2024)} --- two-stage scenario→case derivation with Few-Shot + CoT
and leave-one-out cross-validation over example test cases
(\texttt{temperature=0.1,\ top\_p=1.0,\ max=32768}). - \textbf{Baseline
2 (Milchevski et al., ACL 2025 Industry)} --- intermediate-artifact
workflow: decision table → scenarios → test purposes → Generation Agent
→ Reflection Agent (3 retries each; few-shot disabled as in the
original). - \textbf{Baseline 3 (Bhatia et al., 2024)} ---
conversational: familiarize the model with the full SRS, then chain
per-use-case prompts producing a 6-column Markdown test table
(\texttt{temperature=0.7}); chat memory preserved.

Execution baselines referenced from the thesis: \textbf{AutoUAT}
(scenario-to-code, static Cypress) and \textbf{LLMDroid} (exploratory
coverage agent).

\textbf{Logging.} Every LLM call records prompt/response,
prompt/completion/total tokens, USD cost, and latency to JSONL. Pricing
uses \textbf{placeholder} per-1M rates (flash 0.15/0.60, pro 0.55/2.20,
dola 0.30/1.20 input/output) because official BytePlus rates for these
preview model IDs are not published; \textbf{raw tokens are always
logged, so USD is exactly recomputable} once real rates are known.

\subsubsection{IV-B. Generation results (reproduction on
BytePlus)}\label{iv-b.-generation-results-reproduction-on-byteplus}

\textbf{Full-pipeline resource consumption (measured).}

{\def\LTcaptype{none} % do not increment counter
\begin{longtable}[]{@{}
  >{\raggedright\arraybackslash}p{(\linewidth - 18\tabcolsep) * \real{0.0789}}
  >{\raggedright\arraybackslash}p{(\linewidth - 18\tabcolsep) * \real{0.0789}}
  >{\raggedleft\arraybackslash}p{(\linewidth - 18\tabcolsep) * \real{0.1053}}
  >{\raggedleft\arraybackslash}p{(\linewidth - 18\tabcolsep) * \real{0.1053}}
  >{\raggedleft\arraybackslash}p{(\linewidth - 18\tabcolsep) * \real{0.1053}}
  >{\raggedleft\arraybackslash}p{(\linewidth - 18\tabcolsep) * \real{0.1053}}
  >{\raggedleft\arraybackslash}p{(\linewidth - 18\tabcolsep) * \real{0.1053}}
  >{\raggedleft\arraybackslash}p{(\linewidth - 18\tabcolsep) * \real{0.1053}}
  >{\raggedleft\arraybackslash}p{(\linewidth - 18\tabcolsep) * \real{0.1053}}
  >{\raggedleft\arraybackslash}p{(\linewidth - 18\tabcolsep) * \real{0.1053}}@{}}
\toprule\noalign{}
\begin{minipage}[b]{\linewidth}\raggedright
Case study
\end{minipage} & \begin{minipage}[b]{\linewidth}\raggedright
Model
\end{minipage} & \begin{minipage}[b]{\linewidth}\raggedleft
Calls
\end{minipage} & \begin{minipage}[b]{\linewidth}\raggedleft
Prompt tok
\end{minipage} & \begin{minipage}[b]{\linewidth}\raggedleft
Completion tok
\end{minipage} & \begin{minipage}[b]{\linewidth}\raggedleft
Total tokens
\end{minipage} & \begin{minipage}[b]{\linewidth}\raggedleft
Cost (placeholder)
\end{minipage} & \begin{minipage}[b]{\linewidth}\raggedleft
Wall time
\end{minipage} & \begin{minipage}[b]{\linewidth}\raggedleft
Test cases
\end{minipage} & \begin{minipage}[b]{\linewidth}\raggedleft
Mutations
\end{minipage} \\
\midrule\noalign{}
\endhead
\bottomrule\noalign{}
\endlastfoot
FullTeaching (31pg) & flash & 69 & 1,873,422 & 576,890 &
\textbf{2,450,312} & \textasciitilde\$0.63 & \textasciitilde50 min & 63
& 116 \\
FullTeaching (31pg) & pro & 90 & 2,192,937 & 448,668 &
\textasciitilde{}\textbf{2,641,605}* & \textasciitilde\$2.19* &
\textasciitilde63 min & 63 & 104 \\
4LV sample (2 UC) & flash & 40 & 201,417 & 227,680 & \textbf{429,097} &
\textasciitilde\$0.17 & \textasciitilde28 min & 2 & 3 \\
4LV sample (2 UC) & pro & 10 & 32,527 & 34,247 & \textbf{66,774} &
\textasciitilde\$0.09 & \textasciitilde7 min & 2 & 3 \\
\end{longtable}
}

* The pro FullTeaching total is cumulative with its 4LV run in the
shared usage log; the FullTeaching-only increment is the dominant
component. This is disclosed rather than silently netted out.

\textbf{Per-module token distribution (FullTeaching, flash).} The
generation cost is heavily concentrated in test-case synthesis:

{\def\LTcaptype{none} % do not increment counter
\begin{longtable}[]{@{}lrr@{}}
\toprule\noalign{}
Module & Tokens & Share \\
\midrule\noalign{}
\endhead
\bottomrule\noalign{}
\endlastfoot
test\_case\_generator & 1,640,600 & \textasciitilde67\% \\
mutation\_test\_case\_generator & 455,782 & \textasciitilde19\% \\
screen\_variable\_detector & 166,693 & \textasciitilde7\% \\
screen\_graph\_detector & 120,251 & \textasciitilde5\% \\
business\_process\_detector & 45,579 & \textasciitilde2\% \\
screen\_extraction & 21,407 & \textasciitilde1\% \\
\end{longtable}
}

(Figure F3 visualizes this distribution.)

\textbf{Cross-model finding.} flash and pro produce \textbf{comparable
artifact volumes} (both 63 ISO-style test cases; similar mutation
counts), but pro costs \textbf{\textasciitilde3.5× more} per suite under
placeholder rates while issuing fewer, more concise calls (90 calls at
\textasciitilde2.6M tokens for pro vs 69 calls at 2.45M for flash ---
pro front-loads more tokens per call). For a cost-sensitive, high-volume
regression-suite generation task, flash is the economically preferable
model at comparable output volume; pro's value would need to be
justified by a downstream quality delta not measured here.

\textbf{Baseline re-runs on BytePlus (flash), confirming Option-A wiring
works with zero code change:}

{\def\LTcaptype{none} % do not increment counter
\begin{longtable}[]{@{}
  >{\raggedright\arraybackslash}p{(\linewidth - 8\tabcolsep) * \real{0.1875}}
  >{\raggedright\arraybackslash}p{(\linewidth - 8\tabcolsep) * \real{0.1875}}
  >{\raggedright\arraybackslash}p{(\linewidth - 8\tabcolsep) * \real{0.1875}}
  >{\raggedleft\arraybackslash}p{(\linewidth - 8\tabcolsep) * \real{0.2500}}
  >{\raggedright\arraybackslash}p{(\linewidth - 8\tabcolsep) * \real{0.1875}}@{}}
\toprule\noalign{}
\begin{minipage}[b]{\linewidth}\raggedright
Baseline
\end{minipage} & \begin{minipage}[b]{\linewidth}\raggedright
Method
\end{minipage} & \begin{minipage}[b]{\linewidth}\raggedright
Output
\end{minipage} & \begin{minipage}[b]{\linewidth}\raggedleft
Cost (placeholder)
\end{minipage} & \begin{minipage}[b]{\linewidth}\raggedright
Tokens (in/out)
\end{minipage} \\
\midrule\noalign{}
\endhead
\bottomrule\noalign{}
\endlastfoot
1 (Augusto) & scenario→case (RQ1→RQ2) & 26 test-case files &
\textasciitilde\$1.09 & 140,081 / 414,787 \\
2 (Milchevski) & decision table + gen/reflect specs & decision table +
refined specs & \textasciitilde\$0.28 & 20,943 / 110,129 \\
3 (Bhatia) & conversational per-use-case & 33 test cases &
\textasciitilde\$0.018 & \textasciitilde3,900 / \textasciitilde6,900 \\
\end{longtable}
}

Baseline 1 required the intended manual \textbf{RQ1→RQ2 bridge} (the
model emitted scenarios in a \texttt{**SCENARIO\ N:**} format; copying
the RQ1 output into \texttt{inputTestScenarios.txt} allowed RQ2 title
extraction to succeed).

\subsubsection{IV-C. Execution results (TEAA on live
FullTeaching)}\label{iv-c.-execution-results-teaa-on-live-fullteaching}

Verified programmatically from
\texttt{my\_method\_evaluation\_summary.json}; all aggregates reconcile
exactly.

\textbf{Top-line (all aggregates internally consistent):}

{\def\LTcaptype{none} % do not increment counter
\begin{longtable}[]{@{}
  >{\raggedright\arraybackslash}p{(\linewidth - 4\tabcolsep) * \real{0.3333}}
  >{\raggedright\arraybackslash}p{(\linewidth - 4\tabcolsep) * \real{0.3333}}
  >{\raggedright\arraybackslash}p{(\linewidth - 4\tabcolsep) * \real{0.3333}}@{}}
\toprule\noalign{}
\begin{minipage}[b]{\linewidth}\raggedright
Metric
\end{minipage} & \begin{minipage}[b]{\linewidth}\raggedright
Value
\end{minipage} & \begin{minipage}[b]{\linewidth}\raggedright
Reconciliation
\end{minipage} \\
\midrule\noalign{}
\endhead
\bottomrule\noalign{}
\endlastfoot
Total test cases & \textbf{166} & 70 passed + 96 failed = 166 ✓ \\
Success rate & \textbf{42.17\%} & 70/166 ✓ \\
Valid-action rate & \textbf{81.33\%} & 135 valid + 31 invalid = 166;
135/166 ✓ \\
LLM failures & 35 & 35 + 61 web = 96 failed ✓ \\
Web failures & 61 & --- \\
Total browser actions & \textbf{1,756} & done 528 + input 553 + click
667 + upload 8 ✓ \\
\end{longtable}
}

\textbf{Action-type success (step-level):}

{\def\LTcaptype{none} % do not increment counter
\begin{longtable}[]{@{}llr@{}}
\toprule\noalign{}
Action & Success / Total & Rate \\
\midrule\noalign{}
\endhead
\bottomrule\noalign{}
\endlastfoot
done & 363 / 528 & 68.8\% \\
click & 220 / 667 & 33.0\% \\
input & 142 / 553 & \textbf{25.7\%} \\
upload & 4 / 8 & 50.0\% \\
\textbf{All} & \textbf{729 / 1,756} & \textbf{41.5\%} \\
\end{longtable}
}

\texttt{input} is the weakest action, consistent with the
failure-pattern data (269 input failures) and the frequent
\texttt{no\ such\ element\ {[}id=email/password/...{]}} errors --- the
agent often targets input fields by an \texttt{id} that is not present
in the current DOM state.

\textbf{Selector usage vs.~failures}
(\texttt{my\_method\_evaluation\_failure\_patterns.json}): usage
\texttt{id} 1496 / \texttt{text} 158 / \texttt{xpath} 78 /
\texttt{class} 13; failures \texttt{id} 590 / \texttt{xpath} 49 /
\texttt{text} 48 / \texttt{class} 5. Failures by step type:
\texttt{test\_action} 402, \texttt{navigation} 99, \texttt{path\_action}
11. Dominant runtime errors: \texttt{NoSuchElement} on hallucinated
\texttt{id}s, \texttt{element\ not\ interactable},
\texttt{stale\ element\ reference}, brittle deep XPath chains, and one
hard-coded local upload path that did not exist. (Figure F5: pass/fail
by business process --- BP1 carries 132 of 166 cases; Figure F6:
action-type success.)

\textbf{Static-script baseline evidence.} \texttt{test\_result.xlsx} /
\texttt{test\_result\_updated.xlsx} hold \textbf{169 Cypress-style
results} (failed 132, passed 26, skipped 11), with errors referencing
Cypress, \texttt{loginAndNavigateToCreateCourse()}, and elements
``obscured by the logo header.'' These are the \textbf{static-script
(AutoUAT-style) baseline}, not the TEAA agent, and they evidence the
thesis's ``Static Script Fallacy'' --- high failure from race conditions
and obscured elements, motivating the self-healing agent layer.

\begin{center}\rule{0.5\linewidth}{0.5pt}\end{center}

\subsection{V. Discussion}\label{v.-discussion}

\subsubsection{V-A. Honest reconciliation of execution metrics (the
central integrity
point)}\label{v-a.-honest-reconciliation-of-execution-metrics-the-central-integrity-point}

The thesis reports higher execution figures than the raw artifact
\texttt{my\_method\_evaluation\_summary.json}. These are \textbf{not the
same population}, and per the instruction to \emph{not hide any data},
we lay both out side by side:

{\def\LTcaptype{none} % do not increment counter
\begin{longtable}[]{@{}
  >{\raggedright\arraybackslash}p{(\linewidth - 4\tabcolsep) * \real{0.3333}}
  >{\raggedright\arraybackslash}p{(\linewidth - 4\tabcolsep) * \real{0.3333}}
  >{\raggedright\arraybackslash}p{(\linewidth - 4\tabcolsep) * \real{0.3333}}@{}}
\toprule\noalign{}
\begin{minipage}[b]{\linewidth}\raggedright
Figure
\end{minipage} & \begin{minipage}[b]{\linewidth}\raggedright
Thesis claim
\end{minipage} & \begin{minipage}[b]{\linewidth}\raggedright
Artifact (\texttt{my\_method\_evaluation\_summary.json})
\end{minipage} \\
\midrule\noalign{}
\endhead
\bottomrule\noalign{}
\endlastfoot
Cohort & Normal Flow only, N=166 & All 166 attempted cases (mixed) \\
Pass / pass rate & 148 passed → \textbf{89.20\%} & 70 passed →
\textbf{42.17\%} \\
Valid-action rate & 153/166 → \textbf{92.17\%} & 135/166 →
\textbf{81.33\%} \\
LLM / web failures & 17 / 72 & 35 / 61 \\
Verified BVA Execution Rate (VBER) & \textbf{83.1\%} (64/77 BVA rules) &
not in this file (BVA lives in \texttt{comprehensive\_report.csv}) \\
Boundary Value Coverage (BVC) & \textbf{92.2\%} (71/77) & not in this
file \\
Exception-flow success & thesis Table 4.4 Exception = \textbf{5.9\%}
(3/51) & --- \\
\end{longtable}
}

Four points of caution for anyone citing these numbers: 1. \textbf{The
artifact's 166 ≠ the thesis's ``Normal Flow'' 166.} The artifact's 166
is the \emph{full attempted set} (42.17\% pass). The thesis partitions
results into Normal (N=166, 89.20\%) and Exception (N=51, 5.9\%). The
two ``166''s collide numerically but represent different filtering. One
must not present the artifact's 42.17\% as if it were the 89.20\%
normal-flow figure, nor vice versa. 2. \textbf{92.17\% valid-action}
(thesis) is 153/166 on the \emph{Normal Flow} cohort; the artifact's
comparable field is 81.33\% (135/166) on the \emph{full} set. Both are
internally consistent with their own denominators. 3. \textbf{VBER
83.1\% and BVC 92.2\%} come from the \textbf{FullTeaching BVA benchmark
(N=77 rules)} --- a coverage/ verified-execution axis, not the 166-case
run. 4. \textbf{The ``\textgreater80\% exception validation'' claim is
contradicted by the thesis's own Exception-Flow result (5.9\%).} The
honest reading is that exception handling is a \textbf{documented
limitation} (``Cognitive Ceiling''), not a strength. Any paper built on
this data should present it that way.

Our position: the integrated system reliably executes and verifies
\textbf{normal, happy-path and input-validation flows} at a high rate,
but \textbf{exception/state-driven flows remain hard} --- consistent
with the broader literature (PinATA \textasciitilde60\% verdict
reliability; WebVoyager 59.1\% task success). The 42.2\% all-flows
number is best understood as a \emph{mixed-population} figure dragged
down by the exception cohort and by \texttt{input}-action brittleness
(25.7\% success), not as a refutation of the approach on the flows it
targets.

\subsubsection{V-B. Economic
feasibility}\label{v-b.-economic-feasibility}

Under placeholder pricing, a \textbf{full FullTeaching generation suite
costs \textasciitilde\$0.63 (flash) to \textasciitilde\$2.19 (pro)},
with generation cost concentrated in \texttt{test\_case\_generator}
(\textasciitilde67\%) and \texttt{mutation\_test\_case\_generator}
(\textasciitilde19\%). Because raw tokens are logged, exact USD is
recomputable once official BytePlus rates are published. The economic
argument for the integrated approach rests on \textbf{deterministic
serialization} (eq. 3.20): the expensive LLM-driven discovery and oracle
steps are paid \textbf{once}; verified traces are crystallized into
LLM-free replay scripts (\texttt{ReusableTestRunner}/\texttt{rerun.py})
whose regression cost \texttt{C\_script\ ≪\ C\_LLM} is effectively free.
Thus the one-time generation+discovery spend amortizes across all
subsequent regression cycles --- the key to the thesis's claimed
\textasciitilde88\% cost reduction versus manual testing, which our
token accounting makes independently auditable.

The \textbf{model-selection} implication is concrete: at comparable
artifact volume (63 test cases either way), flash is \textasciitilde3.5×
cheaper than pro, so a cost-optimal deployment uses flash for bulk
generation and reserves a stronger model only where a measured quality
gain justifies it.

\subsubsection{V-C. Threats to validity}\label{v-c.-threats-to-validity}

\begin{enumerate}
\def\labelenumi{\arabic{enumi}.}
\tightlist
\item
  \textbf{4LV is a 2-use-case sample}, not the full RetailOnboardPro SRS
  (10 UC / 8 screens / 75 rules); its results characterize minimal-input
  behavior only and are not a thesis reproduction of RetailOnboardPro.
\item
  \textbf{Placeholder pricing.} All USD figures depend on unofficial
  per-1M rates; only token counts are ground truth.
\item
  \textbf{Cumulative pro token log.} The pro FullTeaching total includes
  its 4LV run in the shared JSONL; the FullTeaching-only figure is the
  dominant increment but is not cleanly isolated.
\item
  \textbf{Execution backend mismatch.} The execution artifacts were
  produced with a Gemini backend. A faithful BytePlus execution re-run
  requires porting \texttt{LLMCaller.py} to the OpenAI-compatible client
  \textbf{and a vision-capable model} for the multi-modal oracle;
  without image support the dual-channel oracle degrades to DOM-only,
  weakening exactly the exception flows already identified as weak.
\item
  \textbf{dola-seed-2-1-turbo} was connectivity-validated but not run
  through a full pipeline; the 3-model comparison is therefore 2-model
  in practice.
\item
  \textbf{NL-test soundness/consistency} (arXiv:2509.19136): some
  execution failures may be false failures from NL ambiguity, and
  re-runs may not be deterministic --- an inherent limitation of
  NL-driven execution.
\end{enumerate}

\begin{center}\rule{0.5\linewidth}{0.5pt}\end{center}

\subsection{VI. Conclusion}\label{vi.-conclusion}

We reproduced the previously published IntelliTest \textbf{generation}
pipeline on a new, provider-agnostic multi-model BytePlus backend with
complete token/cost/latency logging; we \textbf{integrated} it
end-to-end with the TEAA \textbf{execution} agent against a live,
locally-deployed FullTeaching system under test; and we reported a
\textbf{transparent} economic and reliability analysis. The generation
half is prior work and is credited as such throughout. The contributions
of this document are the reproduction (including a real defect fix that
was silently zeroing the pipeline on the new backend), the integration
that closes the generation→execution loop, the first per-model cost
profile for this method (flash ≈ \$0.63, pro ≈ \$2.19 per full
FullTeaching suite; comparable output; pro \textasciitilde3.5×
costlier), and an honest reliability accounting that discloses the
42.2\%-vs-89.2\% population discrepancy and the exception-flow
contradiction rather than reporting only the favourable figure. The
strongest concurrent competitor for the execution phase is \textbf{Yun
Lin's WebTestPilot (FSE 2026)}; a faithful head-to-head against it ---
and measurable improvement of exception-flow reliability via a
vision-capable oracle --- are the clearest next steps.

\begin{center}\rule{0.5\linewidth}{0.5pt}\end{center}

\subsection{References}\label{references}

The complete machine-readable reference set (82 papers, 11-column
schema) is in
\texttt{research\_project/track\_A\_literature\_review/related\_articles.xlsx}
(sheets: Systematic Review {[}82{]}, New Research 2025-26 {[}19{]}, All
References {[}101 rows{]}). Primary anchors cited above:

\begin{itemize}
\tightlist
\item
  Prior work: \emph{A Business Process-Centric Approach for LLM-Driven
  Functional Test Generation} (IEEE Access).
\item
  X. Teoh, \textbf{Y. Lin}, D.-M. Nguyen, R. Ren, W. Zhang, J. S. Dong.
  \emph{WebTestPilot: Agentic End-to-End Web Testing against Natural
  Language Specification by Inferring Oracles with Symbolized GUI
  Elements.} FSE 2026. arXiv:2602.11724.
\item
  \emph{Are Autonomous Web Agents Good Testers? (PinATA / SeeAct-ATA).}
  arXiv:2504.01495.
\item
  C. Augusto et al., \emph{Software System Testing assisted by LLMs: An
  Exploratory Study.} ICTSS 2024.
\item
  D. Milchevski et al., \emph{Multi-Step Generation of Test
  Specifications using LLMs for System-Level Requirements.} ACL 2025
  (Industry).
\item
  S. Bhatia, T. Gandhi, D. Kumar, P. Jalote. \emph{System test case
  design from requirements specifications: insights and challenges of
  using ChatGPT.} 2024.
\item
  J. Wang et al., \emph{Software Testing with LLMs: Survey, Landscape,
  and Vision.} IEEE TSE 2024.
\item
  X. Hou et al., \emph{LLMs for Software Engineering: A Systematic
  Literature Review.} ACM TOSEM 2024.
\item
  H. He et al., \emph{WebVoyager: Building an End-to-End Web Agent with
  LMMs.} ACL 2024.
\item
  J. Yoon, R. Feldt, S. Yoo. \emph{DroidAgent} (intent-driven mobile GUI
  testing).
\item
  Z. Liu et al., \emph{GPTDroid} (functionality-aware mobile GUI
  testing). ICSE 2024.
\item
  Emerging concerns: NL soundness/consistency (arXiv:2509.19136);
  self-healing DOM (arXiv:2603.20358); autonomous repair limits
  (arXiv:2605.01471); Prune4Web DOM pruning (arXiv:2511.21398);
  accessibility-tree prompt injection (arXiv:2507.14799);
  requirement→oracle on real bugs (arXiv:2607.10277).
\item
  Yun Lin group (additional): \emph{Compiling Large Multi-Modal
  Requirement Documents\ldots{}} (ISSTA 2026, arXiv:2602.13723);
  \emph{Generating Project-Specific Test Cases with Requirement
  Validation Intention} (ISSTA 2026); \emph{Generalizing Test Cases for
  Comprehensive Test Scenario Coverage} (FSE 2026, arXiv:2604.21771).
\end{itemize}

\begin{center}\rule{0.5\linewidth}{0.5pt}\end{center}

\subsection{Appendix A --- Reproducibility
artifacts}\label{appendix-a-reproducibility-artifacts}

{\def\LTcaptype{none} % do not increment counter
\begin{longtable}[]{@{}
  >{\raggedright\arraybackslash}p{(\linewidth - 2\tabcolsep) * \real{0.5000}}
  >{\raggedright\arraybackslash}p{(\linewidth - 2\tabcolsep) * \real{0.5000}}@{}}
\toprule\noalign{}
\begin{minipage}[b]{\linewidth}\raggedright
Artifact
\end{minipage} & \begin{minipage}[b]{\linewidth}\raggedright
Path
\end{minipage} \\
\midrule\noalign{}
\endhead
\bottomrule\noalign{}
\endlastfoot
Generation runner + per-model token summaries &
\texttt{research\_project/track\_C\_intellitest\_run/}
(\texttt{run\_intellitest.py}, \texttt{summarize\_usage.py},
\texttt{logs/usage\_*\_summary.json}) \\
4LV PDF renderer &
\texttt{research\_project/track\_C\_intellitest\_run/make\_4lv\_pdf.py} \\
Baseline fidelity audit &
\texttt{research\_project/track\_B\_baselines/REPORT.md} \\
Execution study + architecture→formula mapping &
\texttt{research\_project/track\_D\_execution\_teaa/REPORT.md} \\
Literature review (82 papers) &
\texttt{research\_project/track\_A\_literature\_review/\{systematic\_review.csv,\ related\_articles.xlsx,\ REPORT.md\}} \\
New-research + Yun Lin dossier &
\texttt{research\_project/track\_E\_new\_research/\{new\_papers.csv,\ REPORT.md,\ yun\_lin\_sjtu\_dossier.md\}} \\
Figures (from real logs) &
\texttt{research\_project/track\_F\_plots/figures/F1..F6*.png}
(\texttt{make\_plots.py}) \\
Per-call LLM logs &
\texttt{research\_project/track\_C\_intellitest\_run/logs/usage\_*.jsonl} \\
Shared BytePlus client + logger &
\texttt{research\_project/shared/llm\_client.py} \\
Coordination channels &
\texttt{research\_project/channels/\{MAIN\_LOG.txt,\ writer\_inbox.txt\}} \\
Live SUT & FullTeaching via Docker at
\texttt{https://localhost:5000/} \\
\end{longtable}
}

\subsection{Appendix B --- Figure index}\label{appendix-b-figure-index}

\begin{itemize}
\tightlist
\item
  \textbf{F1} Generation: total tokens per model (FullTeaching latest
  run).
\item
  \textbf{F2} Generation: USD cost per model (placeholder rates, labeled
  as such).
\item
  \textbf{F3} Generation: tokens by pipeline module
  (test\_case\_generator dominates \textasciitilde67\%).
\item
  \textbf{F4} Generation: wall-clock minutes per run.
\item
  \textbf{F5} Execution: pass/fail by business process (overall 42.2\%
  success, N=166; BP1 = 132/166).
\item
  \textbf{F6} Execution: success rate by action type (done
  \textasciitilde69\%, click \textasciitilde33\%, input
  \textasciitilde26\%, upload 50\%).
\end{itemize}

\subsection{Appendix C --- Environment}\label{appendix-c-environment}

Windows; Python 3.11 venv (\texttt{env/}). Installed for this study:
\texttt{matplotlib}, \texttt{openpyxl}, \texttt{fpdf2} (in addition to
the pipeline's \texttt{pydantic}, \texttt{python-dotenv},
\texttt{allpairspy}, \texttt{rapidfuzz}, \texttt{pypdf},
\texttt{google-genai}, \texttt{openai}). Docker 29.7.2 for the SUT (git
required \texttt{core.longpaths=true} on Windows to check out
FullTeaching's deep Java test paths).

\end{document}
